\section{Conclusion}
\label{sec:conclusion}
Cet essai de 115~minutes sur une micro-cogénération au gaz naturel
(6~kW électriques nominaux) nous permet de comprendre les principes de la cogénération, mais
l'installation n'atteint pas les performances annoncées. Calculé sur les
énergies cumulées, le rendement global est de 65,5~\% (24,5~\%
électrique et 41,0~\% thermique), contre 85~\% pour le constructeur. La
puissance thermique récupérée (9,2~kW) est inférieure de 21~\% à la
valeur nominale, et le débit estimé (environ 14~L/min au lieu de
34,1~L/min).

À performances mesurées, l'économie d'énergie primaire est négative
($-8{,}4$~\%, contre $+14{,}7$~\% avec les valeurs du constructeur,
section~\ref{sec:pes}). L'intérêt d'une telle unité dépend donc d'un
fonctionnement proche du point nominal, que la machine ne sait pas atteindre dans son état actuel, d'un besoin de chaleur régulier et
d'un rapport favorable entre les prix de l'électricité et du gaz.
 
Ilserait intéressant de réaliser un nouvel essai avec le débit
d'eau mesuré afin de connaitre la puissance thermique réellement récupérée.
